\origpage{26}{}
\textbf{Proof.} Choose $n \ge 0$ as in \del{(2.5)}\fix{(2.6)} with $k$ replaced by $2q$.\fixnote{fix:6.5-n}{The reference meant is (2.6). Result (2.5) is about separations of order exactly $k$, but the separations $(A_i, B_i)$ below have order $q_1 + q_2$, which may be less than $2q$ (for instance $q_1 = q_2 = 1$ with $q = 2$). Result (2.6) allows order $\le k$.} Choose $\lambda'$ so that (6.4) holds, with $p, q, \lambda$ replaced by $p+2q, q, \lambda'$. Let $\lambda = n(2\lambda'+1)$. We claim that $\lambda$ satisfies (6.5). For let $\Gamma, F_1^*, F_2^*, H_1, H_2, \mu$ be a presentation of $G$ in the disc $\Sigma$, of depth $\le q$, and let $L$ be an exhaustive vital $p$-linkage. Let the two spaces not included in $F_1^* \cup F_2^*$ be $s_1^*$ and $s_2^*$. Let $r_1, \ldots, r_t$ be all the regions of $\Gamma$ in $\Sigma$ which are incident both with a vertex in $F_1^*$ and with a vertex in $F_2^*$, numbered in order, so that for $1 \le i < t, r_i$ lies between $s_1^*$ and $r_{i+1}$, in the natural sense. For $1 \le i \le t$, choose $a_i \in V(\Gamma) \cap F_1^*$ and $b_i \in V(\Gamma) \cap F_2^*$ so that $r_i$ is incident with $a_i$ and $b_i$; and let $F_i$ be a good line with ends $a_i, b_i$ and with interior in $r_i$. Let $J_i$ and $K_i$ be the two lines in $bd(\Sigma)$ with ends $a_i$ and $b_i$, where $s_1^* \subseteq J_i$ and $s_2^* \subseteq K_i$. Then $J_1 \subseteq J_2 \subseteq \ldots \subseteq J_t$ and $K_t \subseteq K_{t-1} \subseteq \ldots \subseteq K_1$. Let $A_i$ be the union of $\mu(J_i)$ and $\Gamma \cap \Delta$, where $\Delta \subseteq \Sigma$ is the disc bounded by $F_i \cup J_i$; and define $B_i$ similarly using $K_i$ instead of $J_i$. Then $(A_i, B_i)$ is a separation of $G$, and $A_i \cap B_i = \mu(a_i) \cup \mu(b_i)$, as is easily seen. Since $\mu(a_i) \cap \mu(b_i) = \emptyset$, it follows that $|V(A_i \cap B_i)| = q_1 + q_2$, where $|\mu(a)| = q_j$ for all $a \in F_j^* \cap V(\Gamma)$ ($j = 1, 2$). From (6.2), for $1 \le i < t$, there are $q_j$ mutually vertex-disjoint paths of $H_j \cap B_i \cap A_{i+1}$ between $V(A_i \cap B_i)$ and $V(A_{i+1} \cap B_{i+1})$ for $j = 1$ and $2$, and since $H_1$ and $H_2$ are disjoint, it follows that there are $q_1 + q_2$ mutually vertex-disjoint paths of $(H_1 \cup H_2) \cap B_i \cap A_{i+1}$ between $V(A_i \cap B_i)$ and $V(A_{i+1} \cap B_{i+1})$. But for $1 \le i < t, \Gamma \cap B_i \cap A_{i+1}$ has an edge since $r_i \ne r_{i+1}$ and it belongs to $L$ since $L$ is exhaustive, and so $L \cap B_i \cap A_{i+1} \not\subseteq H_1 \cup H_2$. From \del{(2.5)}\fix{(2.6)}, it follows that $t \le n$.\fixnote{fix:6.5-t}{As above, the result meant is (2.6), with which $n$ was chosen.}

For $1 \le i < t$, let $P_i$ be the $\Sigma$-jump with $\bar{r}_i \cap \bar{r}_{i+1} = U(P_i)$, and let $\Gamma' = P_1 \cup \ldots \cup P_{t-1}$. Let $R_1, \ldots, R_t$ be the regions of $\Gamma'$ in $\Sigma$, where $r_i \subseteq R_i$ ($1 \le i \le t$). Now let us fix $i$ with $1 \le i \le t$. We claim that there are at most $2\lambda'$ $\Sigma$-jumps $J$ with $U(J) \subseteq R_i$. For $j = 1, 2$, let $I_j \subseteq F_j^*$ be the support $\bar{r}_i \cap F_j^*$; then $\bar{r}_{\delm{j}\fix{i}}$\fixnote{fix:6.5-disc}{The disc bounded by $I_1 \cup U(P_{i-1}) \cup I_2 \cup U(P_i)$ is $\bar{r}_i$, for the fixed $i$. The index $j$ runs over the two supports.} is a closed disc bounded by $I_1 \cup U(P_{i-1}) \cup I_2 \cup U(P_i)$ (replacing $U(P_{i-1})$ by $s_1^*$ if $i = 1$, and replacing $U(P_i)$ by $s_2^*$ if $i = t$). For $j = 1, 2$, let $\Gamma_j$ be the drawing in $\Sigma$ formed by all the vertices in $V(\Gamma) \cap I_j$ and all the $\Sigma$-jumps $J$ with $U(J) \subseteq R_i \cup I_1 \cup I_2$ with both ends in $I_j$. We claim that $\Gamma_j$ includes at most $\lambda'$ $\Sigma$-jumps. For if $|I_j| = 1$ this is trivial, since $\Gamma_j$ includes no $\Sigma$-jumps. If $|I_j| \ne 1$ then $\Gamma_j, I_j, \mu(I_j)$, and the restriction of $\mu$ to $I_j$, is a presentation of $\Gamma_j \cup \mu(I_j)$ in a disc with depth $\le q$ and defect 1; and $L \cap (\Gamma_j \cup \mu(I_j))$ is an exhaustive vital $(p+2q)$-linkage, and so by (6.4) it has at most $\lambda'$ $\Sigma$-jumps. This proves our claim that $\Gamma_1$ and $\Gamma_2$ both include at most $\lambda'$ $\Sigma$-jumps. Since every $\Sigma$-jump $J$ with $U(J) \subseteq R_i \cup I_1 \cup I_2$ either has both ends in $I_1$ or has both ends in $I_2$ (because $r_i$ is the only region of $\Gamma$ included in $R_i$ incident with both $F_1^*$ and $F_2^*$) it follows that $R_i$ includes at most $2\lambda'$ $\Sigma$-jumps. Counting $P_1, \ldots, P_{t-1}$, and using the fact that $t \le n$, we deduce that there are at most $2\lambda't + t - 1 \le \lambda$ $\Sigma$-jumps altogether, as required. $\square$
